\section{Annexes}
This last chapter is the one to open with an exam paper in front of you. It carries
three pieces of reference material --- the lemme de classe monotone, the table of
lois discrètes and the table of lois à densité --- and adds to them a summary the
course gives nowhere, a single page on the modes of convergence.\\
Two things are worth saying before the tables. The first is that the lemme de
classe monotone is not decoration: chapter 2's théorème d'unicité for mesures is
proved by building a classe monotone and invoking the lemma, and without the lemma
that proof has a hole in the middle. It is stated and proved here in full,
hypotheses included.\\
The second is that the two tables of lois are \emph{not} a bare list of
definitions. Such a list gives a name, the parameters, the notation, the support
and the loi; the tables below add the espérance, the variance and the
fonction caractéristique on both sides, and --- for the lois discrètes only --- the
fonction génératrice, from which their fonction caractéristique is read off
directly. Those added columns are computed in chapters 1, 6 and 7, and they are
the reason to have the tables at all: the first five columns alone tell you nothing
you could not read off a definition.

\subsection{Classe monotone}
On se donne un ensemble $\Omega$; all the families considered in this subsection
are families of subsets of $\Omega$.

\begin{definition}[Classe monotone]
  Une famille $\mathcal{M}$ de sous-ensembles de $\Omega$ est appelée classe
  monotone sur $\Omega$ si elle vérifie les propriétés suivantes :
  \begin{itemize}
    \item $\Omega \in \mathcal{M}$ ;
    \item pour tous $A, B \in \mathcal{M}$ tels que $A \subr B$,
      $B \setminus A \in \mathcal{M}$ ;
    \item pour tous $A_{1}, A_{2}, \dots \in \mathcal{M}$ tels que
      $A_{n} \uparrow A$, $A \in \mathcal{M}$.
  \end{itemize}
\end{definition}

\begin{remark}[Why the first condition is needed]
  The first condition is $\Omega \in \mathcal{M}$, membership in the family
  $\mathcal{M}$ itself, and it has to be: the second condition produces complements
  only by subtracting from a set already in the family, so without $\Omega$ in
  $\mathcal{M}$ the family need not contain a single complement and the proposition
  below would be false.
\end{remark}

Il est facile de vérifier que toute tribu est une classe monotone: a tribu contains
$\Omega$; it is stable under proper difference, since
$B \setminus A = B \intersect A^{c}$; and it is stable under increasing limits,
since an increasing limit is a countable union. L'inverse n'est pas vrai en
général --- a classe monotone need not be stable under intersection, and then it is
not a tribu. Une condition nécessaire et suffisante pour qu'une classe monotone soit
une tribu est qu'elle soit stable par intersection finie: the gap between the two
notions is exactly that stability, which is the content of the next statement.

\begin{prop}[Classe monotone qui est une tribu]
  Une classe monotone $\mathcal{M}$ est une tribu si et seulement si elle est stable
  par intersection finie.
\end{prop}

\begin{proof}
  If $\mathcal{M}$ is a tribu then it is stable par intersection finie, since
  $A \intersect B = \left(A^{c} \union B^{c}\right)^{c}$ and a tribu is stable under
  complement and under countable --- hence finite --- union.\\
  Conversely, suppose the classe monotone $\mathcal{M}$ is stable par intersection
  finie. It contains $\Omega$ by the first axiom. It is stable under
  complement: for $A \in \mathcal{M}$ one has $A \subr \Omega$ and both sets lie in
  $\mathcal{M}$, so $A^{c} = \Omega \setminus A \in \mathcal{M}$ by the second
  axiom. Par passage au complémentaire, on obtient donc la stabilité par réunion
  finie, since
  \[
    A \union B = \left(A^{c} \intersect B^{c}\right)^{c},
  \]
  the intersection on the right being in $\mathcal{M}$ by hypothesis. There remains
  countable union. Pour montrer la stabilité par réunion dénombrable, on considère
  une suite $(A_{n})_{n \geq 1}$ dans $\mathcal{M}$. Put
  $B_{n} = \bigcup_{k=1}^{n} A_{k}$: each $B_{n}$ lies in $\mathcal{M}$ by stability
  under finite union. Alors la suite $(B_{n})_{n \geq 1}$ est croissante, si bien
  que sa limite
  \[
    \bigcup_{n \geq 1} A_{n} = \lim_{n \to +\infty} \uparrow B_{n}
  \]
  appartient à $\mathcal{M}$, by the third axiom. So $\mathcal{M}$ contains
  $\Omega$ and is stable under complement and countable union: it is a tribu.
\end{proof}

\begin{remark}
  The finite-union step is where stability under intersection is spent, and it is
  spent only there. That is the shape of every application: one wants a countable
  union, the third axiom delivers increasing limits only, and the gap is closed by
  turning an arbitrary countable union into an increasing one --- which costs finite
  unions, which cost complements and finite intersections.
\end{remark}

Comme pour les tribus, toute intersection de classes monotones sur $\Omega$ est une
classe monotone: each of the three conditions is stable under intersecting the
families, because each is a condition of the form ``these sets belong to the
family''. Since $\mathcal{P}(\Omega)$ is itself a classe monotone containing any
collection whatsoever, the intersection below is taken over a non-empty family, ce
qui permet de définir la notion de classe monotone engendrée.

\begin{definition}[Classe monotone engendrée]
  Soit $\mathcal{C} \subr \mathcal{P}(\Omega)$ une collection d'ensembles.
  L'intersection de toutes les classes monotones contenant $\mathcal{C}$ est une
  classe monotone, appelée la \textbf{classe monotone engendrée}
  (\textit{generated monotone class}) par $\mathcal{C}$. On la note
  $\mathcal{M}(\mathcal{C})$.
\end{definition}

The classe monotone engendrée is contained in the tribu engendrée, since a tribu is
a classe monotone. La classe monotone engendrée n'est pas la tribu engendrée en
général. C'est le cas dès que la collection $\mathcal{C}$ est stable par
intersection finie. C'est le \textbf{lemme de classe monotone} (\textit{monotone
class lemma}).

\begin{lemma}[Lemme de classe monotone]
  Soit $\mathcal{C} \subr \mathcal{P}(\Omega)$ une collection d'ensembles stable par
  intersection finie. Alors
  \[
    \mathcal{M}(\mathcal{C}) = \sigma(\mathcal{C}).
  \]
\end{lemma}

\begin{proof}
  Puisque toute tribu est une classe monotone --- $\sigma(\mathcal{C})$ is one, and
  it contains $\mathcal{C}$ ---, on a nécessairement
  $\mathcal{M}(\mathcal{C}) \subr \sigma(\mathcal{C})$.\\
  Pour montrer l'autre inclusion, il suffit de montrer que $\mathcal{M}(\mathcal{C})$
  est une tribu --- it then contains $\sigma(\mathcal{C})$, the smallest tribu
  containing $\mathcal{C}$ ---, c'est-à-dire, d'après la proposition, qu'elle est
  stable par intersection finie: being a tribu amounts to that for a classe
  monotone, and that is what is proved, in two steps.\\
  \emph{First step: one set in $\mathcal{C}$.} Soit $A \in \mathcal{C}$. On
  définit :
  \[
    \mathcal{M} = \{B \in \mathcal{M}(\mathcal{C}) :
      A \intersect B \in \mathcal{M}(\mathcal{C})\}.
  \]
  Comme $\mathcal{C}$ est stable par intersection finie --- $A \intersect B$ lies
  in $\mathcal{C} \subr \mathcal{M}(\mathcal{C})$ for every $B \in \mathcal{C}$
  ---, on a $\mathcal{C} \subr \mathcal{M}$. On vérifie par ailleurs que
  $\mathcal{M}$ est une classe monotone :
  \begin{itemize}
    \item $\Omega \in \mathcal{M}$, because $A \intersect \Omega = A \in
      \mathcal{C} \subr \mathcal{M}(\mathcal{C})$ ;
    \item pour tous $A', B' \in \mathcal{M}$ tels que $A' \subr B'$,
      \[
        A \intersect (B' \setminus A')
          = (A \intersect B') \setminus (A \intersect A')
          \in \mathcal{M}(\mathcal{C}),
      \]
      the two sets on the right being in $\mathcal{M}(\mathcal{C})$ by definition of
      $\mathcal{M}$ and nested because $A' \subr B'$, donc
      $B' \setminus A' \in \mathcal{M}$ ;
    \item pour tous $A'_{1}, A'_{2}, \dots \in \mathcal{M}$ tels que
      $A'_{n} \uparrow A'$, one has $A \intersect A'_{n} \uparrow A \intersect A'$,
      an increasing limit of sets of $\mathcal{M}(\mathcal{C})$, whence
      $A \intersect A' \in \mathcal{M}(\mathcal{C})$, de sorte que
      $A' \in \mathcal{M}$.
  \end{itemize}
  Comme $\mathcal{M}$ --- a classe monotone --- contient $\mathcal{C}$, on a
  $\mathcal{M}(\mathcal{C}) \subr \mathcal{M}$. Unwinding the definition of
  $\mathcal{M}$, on a donc montré que :
  \[
    \forall A \in \mathcal{C}, \ \forall B \in \mathcal{M}(\mathcal{C}), \qquad
      A \intersect B \in \mathcal{M}(\mathcal{C}).
  \]
  \emph{Second step: both sets in $\mathcal{M}(\mathcal{C})$.} Pour conclure, on
  fixe $B \in \mathcal{M}(\mathcal{C})$ et on définit :
  \[
    \mathcal{M}' = \{A \in \mathcal{M}(\mathcal{C}) :
      A \intersect B \in \mathcal{M}(\mathcal{C})\}.
  \]
  On sait que $\mathcal{C} \subr \mathcal{M}'$: it is exactly the conclusion of the
  first step. Par les mêmes arguments que précédemment --- the same three
  verifications as above, with $B$ now playing the role $A$ played there ---, on
  montre que $\mathcal{M}'$ est une classe monotone, si bien que
  $\mathcal{M}(\mathcal{C}) \subr \mathcal{M}'$. On a donc :
  \[
    \forall A, B \in \mathcal{M}(\mathcal{C}), \qquad
      A \intersect B \in \mathcal{M}(\mathcal{C}).
  \]
  Ce qui prouve que $\mathcal{M}(\mathcal{C})$ est stable par intersection finie.
  Being a classe monotone, it is therefore a tribu, which completes the proof.
\end{proof}

\begin{remark}[Why the two steps]
  The asymmetry is the whole proof. The hypothesis on $\mathcal{C}$ gives
  intersections of two sets \emph{of $\mathcal{C}$}; the first step upgrades one of
  the two arguments to an arbitrary element of $\mathcal{M}(\mathcal{C})$, and the
  second step upgrades the other, using the first step's conclusion as the seed
  $\mathcal{C} \subr \mathcal{M}'$. Dropping stability of $\mathcal{C}$ under finite
  intersection kills the first step at its first line, and the lemma is false
  without it --- as is chapter 2's théorème d'unicité, which is why that theorem
  carries the same hypothesis, alongside its own second condition that $\Omega$ be
  covered by countably many elements of $\mathcal{C}$ of finite mesure. Both are
  load-bearing there: the first buys this lemma, the second is what lets the
  equality be passed to the limit along an increasing exhaustion of $\Omega$.
\end{remark}

\subsection{Lois usuelles}
The two tables below list the usual lois with columns added --- four to the table
of lois discrètes, three to the table of lois à densité. Name, parameters,
notation, support and loi define each entry; the espérance, the variance and the
transforms are computed elsewhere in these notes and are gathered here. Read the
notation column first: it fixes which meaning each of the overloaded symbols
carries, and the closing note of this chapter lists the
overloads.

\subsubsection{Lois discrètes}

\begin{center}
\footnotesize
\setlength{\tabcolsep}{4pt}
\begin{tabular}{llllccc}
  \toprule
  Nom & Paramètre(s) & Notation & Support & Loi & $\expec(X)$ & $\Var(X)$ \\
  \midrule
  Bernoulli & $p \in [0,1]$ & $\mathcal{B}(p)$ & $\{0,1\}$
    & $\begin{cases} p & k = 1\\ 1-p & k = 0\end{cases}$ & $p$ & $p(1-p)$ \\[6pt]
  Binomiale & $n \in \mathbb{N}$, $p \in [0,1]$ & $\mathcal{B}(n,p)$
    & $\{0,1,\dots,n\}$ & $\binom{n}{k}p^{k}(1-p)^{n-k}$ & $np$ & $np(1-p)$ \\[3pt]
  Poisson & $\lambda > 0$ & $\mathcal{P}(\lambda)$ & $\mathbb{N}$
    & $e^{-\lambda}\dfrac{\lambda^{k}}{k!}$ & $\lambda$ & $\lambda$ \\[5pt]
  Géométrique & $p \in\, ]0,1]$ & $\mathcal{G}(p)$ & $\mathbb{N}^{*}$
    & $(1-p)^{k-1}p$ & $\dfrac{1}{p}$ & $\dfrac{1-p}{p^{2}}$ \\[5pt]
  Uniforme & $n \in \mathbb{N}^{*}$ & $\mathcal{U}(1,\dots,n)$ & $\{1,\dots,n\}$
    & $\dfrac{1}{n}$ & $\dfrac{n+1}{2}$ & $\dfrac{n^{2}-1}{12}$ \\
  \bottomrule
\end{tabular}
\end{center}

\begin{center}
\footnotesize
\setlength{\tabcolsep}{6pt}
\begin{tabular}{llcc}
  \toprule
  Nom & Notation & $G_{X}(t) = \expec\!\left(t^{X}\right)$
    & $\varphi_{X}(t) = G_{X}\!\left(e^{it}\right)$ \\
  \midrule
  Bernoulli & $\mathcal{B}(p)$ & $1-p+pt$ & $1-p+pe^{it}$ \\[3pt]
  Binomiale & $\mathcal{B}(n,p)$ & $(1-p+pt)^{n}$ & $\left(1-p+pe^{it}\right)^{n}$
    \\[3pt]
  Poisson & $\mathcal{P}(\lambda)$ & $e^{\lambda(t-1)}$
    & $e^{\lambda\left(e^{it}-1\right)}$ \\[3pt]
  Géométrique & $\mathcal{G}(p)$ & $\dfrac{pt}{1-(1-p)t}$
    & $\dfrac{pe^{it}}{1-(1-p)e^{it}}$ \\[5pt]
  Uniforme & $\mathcal{U}(1,\dots,n)$ & $\dfrac{t}{n}\dfrac{1-t^{n}}{1-t}$
    & $\dfrac{1}{n}e^{it}\dfrac{e^{int}-1}{e^{it}-1}$ \\
  \bottomrule
\end{tabular}
\end{center}

The first table gives the definitions plus the two moment columns; the second carries
the fonction génératrice computed in chapter 1 and, beside it, the fonction
caractéristique of chapter 6. The two are the same object: for a variable with
values in $\mathbb{N}$,
\[
  \varphi_{X}(t) = \expec\!\left(e^{itX}\right)
    = \expec\!\left(\left(e^{it}\right)^{X}\right) = G_{X}\!\left(e^{it}\right),
\]
so the right-hand column is the middle one evaluated on the unit circle and nothing
new has to be memorised. The entry for the loi uniforme is written for $t \neq 1$
and, in the characteristic form, for $t \notin 2\pi\zee$; at the excluded points
every term of the sum is $1$ and the value is $1$.

\begin{remark}[The geometric support is a convention, and it moves]
  The table supports the loi géométrique on $\mathbb{N}^{*}$: $X$ counts the
  trials up to and including the first success, so $\expec(X) = 1/p$. Many
  sources, and some exercises, instead use $\tau = X - 1$, the number of failures
  \emph{before} the
  first success, supported on $\mathbb{N}$, with $\expec(\tau) = (1-p)/p$ and the
  same variance $(1-p)/p^{2}$. The table above follows the course's convention.
  Check which of the two a question means before quoting a mean.
\end{remark}

\subsubsection{Lois à densité}

\begin{center}
\footnotesize
\setlength{\tabcolsep}{4pt}
\begin{tabular}{lllllcc}
  \toprule
  Nom & Paramètre(s) & Notation & Support & Densité & $\expec(X)$ & $\Var(X)$ \\
  \midrule
  Exponentielle & $\lambda > 0$ & $\mathcal{E}(\lambda)$ & $\reals_{+}$
    & $\lambda e^{-\lambda x}\indic_{\reals_{+}}(x)$
    & $\dfrac{1}{\lambda}$ & $\dfrac{1}{\lambda^{2}}$ \\[5pt]
  Gamma & $a, \lambda > 0$ & $\Gamma(a,\lambda)$ & $\reals_{+}$
    & $\dfrac{\lambda^{a}}{\Gamma(a)}x^{a-1}e^{-\lambda x}\indic_{\reals_{+}}(x)$
    & $\dfrac{a}{\lambda}$ & $\dfrac{a}{\lambda^{2}}$ \\[7pt]
  Uniforme & $a < b$ & $\mathcal{U}([a,b])$ & $[a,b]$
    & $\dfrac{1}{b-a}\indic_{[a,b]}(x)$
    & $\dfrac{a+b}{2}$ & $\dfrac{(b-a)^{2}}{12}$ \\[7pt]
  Beta & $a, b > 0$ & $\mathcal{B}(a,b)$ & $[0,1]$
    & $\dfrac{1}{B(a,b)}x^{a-1}(1-x)^{b-1}\indic_{[0,1]}(x)$
    & $\dfrac{a}{a+b}$ & $\dfrac{ab}{(a+b)^{2}(a+b+1)}$ \\[7pt]
  Normale & $\mu \in \reals$, $\sigma^{2} > 0$ & $\mathcal{N}(\mu,\sigma^{2})$
    & $\reals$
    & $\dfrac{1}{\sqrt{2\pi}\,\sigma}\,e^{-\frac{(x-\mu)^{2}}{2\sigma^{2}}}$
    & $\mu$ & $\sigma^{2}$ \\
  \bottomrule
\end{tabular}
\end{center}

Each support is the set on which the densité's indicator does not vanish; the
densité of the loi normale carries no indicator, and its support is all of
$\reals$. The fonctions caractéristiques, computed in chapter 6, are:

\begin{center}
\footnotesize
\setlength{\tabcolsep}{6pt}
\begin{tabular}{llc}
  \toprule
  Nom & Notation & $\varphi_{X}(t) = \expec\!\left(e^{itX}\right)$ \\
  \midrule
  Exponentielle & $\mathcal{E}(\lambda)$ & $\dfrac{\lambda}{\lambda - it}$ \\[5pt]
  Gamma & $\Gamma(a,\lambda)$
    & $\left(\dfrac{\lambda}{\lambda - it}\right)^{a}$ \\[7pt]
  Uniforme & $\mathcal{U}([a,b])$
    & $\dfrac{e^{itb} - e^{ita}}{it\,(b-a)}$ for $t \neq 0$, and $1$ at $t = 0$
    \\[7pt]
  Beta & $\mathcal{B}(a,b)$
    & $\displaystyle\sum_{k \geq 0}
      \left(\prod_{r=0}^{k-1}\frac{a+r}{a+b+r}\right)\frac{(it)^{k}}{k!}$ \\[9pt]
  Normale & $\mathcal{N}(\mu,\sigma^{2})$
    & $e^{i\mu t}e^{-\sigma^{2}t^{2}/2}$ \\
  \bottomrule
\end{tabular}
\end{center}

\begin{remark}[The loi normale row is an addition]
  The loi normale is developed with the vecteurs gaussiens in chapter 7 rather
  than alongside the lois exponentielle, Gamma, uniforme and Beta. The row above is
  added anyway --- its fonction caractéristique is chapter 6's, and the reason for
  wanting it in a reference table at all is chapter 7's Gaussian development ---
  and it is worth having here, because it is the one loi an exam is certain to use
  and the one row a revision sheet built from those four lois alone would lack.
\end{remark}

\begin{remark}[Which entries are computed, and where]
  The moments of the lois exponentielle and Gamma come from the fonction
  caractéristique by the moment expansion of chapter 6, and the Gamma line at
  $a = 1$ reproduces the Exponentielle line as a check. The moments of the loi
  uniforme are chapter 6's unit-interval
  computation $\expec = 1/2$, $\Var = 1/12$ transported by the affine map
  $x \mapsto a + (b-a)x$, which shifts the mean to $(a+b)/2$ and scales the variance
  by $(b-a)^{2}$. The moments of the loi Beta come from the moment formula
  $\expec(X^{k}) = \prod_{r=0}^{k-1}\frac{a+r}{a+b+r}$ of the same chapter, at
  $k = 1$ and $k = 2$. The Normale row's fonction caractéristique is chapter 6's;
  its mean and variance are its parameters, as chapter 7 defines them.
\end{remark}

\begin{remark}[The beta fonction caractéristique has no elementary closed form]
  Four of the five lois above have a fonction caractéristique in closed form; the
  loi Beta does not, and the series recorded in the table is the honest answer
  rather than an omission. It converges everywhere --- the support is
  bounded, so every moment exists and the partial sums are dominated by $e^{|t|}$
  times the densité --- and it is the confluent hypergeometric function
  $M(a,\,a+b,\,it)$. On an exam a question about the loi Beta asks for a
  moment, not for $\varphi_{X}$.
\end{remark}

\subsection{Notions de convergence}
The course defines two modes of convergence for variables aléatoires and names a
third only to say that it is not treated. They are collected here with the
implications between them, because the two theorems of chapter 8 deliver different
ones and saying which is half of a correct answer.\\
La notion de convergence la plus forte --- the strongest of the three --- est une
convergence simple, valable pour presque toute réalisation $\omega$: the numbers
$X_{n}(\omega)$ converge, and the set of $\omega$ at which
$X_{n}(\omega) \to X(\omega)$ has probability $1$. On parle de convergence
presque sûre. On utilise la notation
$X_{n} \cvas X$. It composes freely with continuous maps, as convergence en loi also does.\\
La convergence en loi forgets the univers entirely: elle ne concerne que les lois
des variables aléatoires impliquées, et non leur réalisation. They are compared
through the fonctions de répartition, and
$F_{n}(x) \to F(x)$ is required at the continuity points of $F$ alone. On utilise
la notation $X_{n} \cvL X$; comme cette notion de convergence s'applique aux lois,
on pourra également l'écrire comme tel, par exemple
$X_{n} \cvL \mathcal{N}(0,1)$. The restriction to continuity
points is the content of the definition and not a technicality: without
it the deterministic sequence $1/n$ would fail to converge en loi to $0$.\\
La convergence en probabilité requires $\prob(|X_{n} - X| > a) \to 0$ for every
$a > 0$, and is written $X_{n} \cvP X$.
\emph{The course does not define it.} D'autres notions de convergence existent,
comme la convergence en probabilité, mais ne sont pas abordées dans ce cours: the
course names it only to say so, and proves nothing about it. It
is carried in these notes as an addition, because it sits between the two modes the
course does define, and a problem on this syllabus may well supply its definition
inline before asking for an implication --- which is itself a sign that the course
does not supply it.

\begin{figure}[h!]
  \centering
  \begin{tikzpicture}[>=stealth, node distance=12mm and 26mm]
    \node[draw, very thick, rounded corners, inner sep=6pt] (as)
      {$X_{n} \cvas X$};
    \node[draw, dashed, rounded corners, inner sep=6pt, right=of as] (p)
      {$X_{n} \cvP X$};
    \node[draw, very thick, rounded corners, inner sep=6pt, right=of p] (l)
      {$X_{n} \cvL X$};
    \draw[->, very thick] (as) -- node[above,font=\footnotesize] {conv.\ dominée} (p);
    \draw[->, dashed] (p) -- node[above,font=\footnotesize] {not proved here} (l);
    \draw[->, very thick] (as.south) .. controls +(0,-1.2) and +(0,-1.2) ..
      node[below,font=\footnotesize]
      {conv.\ dominée on bounded continuous tests} (l.south);
    \node[below=18mm of p, align=center, font=\footnotesize]
      {the dashed box and the dashed arrow are the addition ---\\
       the course defines only the two solid modes};
  \end{tikzpicture}
\end{figure}

\begin{remark}[No arrow reverses]
  Convergence presque sûre implies convergence en probabilité --- proved in
  chapter 8 --- and convergence en probabilité implies convergence en loi,
  an implication carried here without proof; neither converse holds. An indicator
  window of width $1/n$ sliding cyclically around $[0,1]$ has $\prob(|X_{n}| > a)
  = 1/n \to 0$, so $X_{n} \cvP 0$, while every $\omega$ is hit infinitely often
  and $X_{n}(\omega)$
  converges for no $\omega$ at all. And with $X$ of loi normale centrée réduite and
  $X_{n} = -X$ for every $n$, the lois are equal so $X_{n} \cvL X$ holds trivially,
  while $|X_{n} - X| = 2|X|$ never gets small. The one partial converse worth
  knowing belongs to the addition and is stated, not proved, in these notes:
  convergence en loi \emph{to a constant} does imply convergence en probabilité to
  that constant.
\end{remark}

The practical criterion for convergence en loi is not the definition but the
fonction caractéristique, and it comes from two results used together. Chapter 6
proves that the fonction caractéristique determines the loi:
$\prob_{X} = \prob_{Y}$ if and only if $\varphi_{X} = \varphi_{Y}$. Chapter 8's
theorem on the three characterisations of convergence en loi then upgrades that
from lois to sequences of lois.

\begin{theorem}[Trois caractérisations de la convergence en loi]
  Soit $X_{1}, X_{2}, \dots$ une suite de variables aléatoires réelles, de
  fonctions caractéristiques respectives $\varphi_{1}, \varphi_{2}, \dots$, et
  $X$ une variable aléatoire réelle, de fonction caractéristique $\varphi$. Les
  propositions suivantes sont équivalentes :
  \begin{enumerate}
    \item $X_{n} \cvL X$ ;
    \item $\expec(\phi(X_{n})) \to \expec(\phi(X))$ pour toute fonction
      \emph{continue bornée} $\phi$ ;
    \item $\varphi_{n}(t) \to \varphi(t)$ pour tout $t \in \reals$.
  \end{enumerate}
\end{theorem}

The third characterisation is the one that gets used: to prove a convergence en
loi, compute the fonction caractéristique of the $n$-th term, take its
pointwise limit, and recognise the result. That is exactly how the théorème
central limite is proved.\\
Two hypotheses inside the criterion are worth keeping in view. The fonctions test
of the second characterisation must be bounded \emph{and} continuous ---
boundedness alone is what makes convergence dominée available, and continuity
alone is what lets convergence presque sûre be pushed through. And the pointwise
convergence of the third is convergence at every real $t$, not on a neighbourhood
of $0$ and not uniformly.\\
Which mode does each theorem deliver? The two results of chapter 8 answer
differently, and confusing them is the standard error.

\begin{theorem}[Loi forte des grands nombres]
  Pour toute suite $X_{1}, X_{2}, \dots$ de variables aléatoires réelles
  i.i.d.\ \emph{ayant une espérance},
  \[
    \frac{1}{n}\sum_{i=1}^{n} X_{i} \cvas \expec(X_{1}).
  \]
\end{theorem}

\begin{theorem}[Théorème central limite]
  Pour toute suite $X_{1}, X_{2}, \dots$ de variables aléatoires réelles
  i.i.d.\ d'espérance finie $\mu$ et de variance finie $\sigma^{2}$,
  \[
    \sqrt{n}\left(\frac{1}{n}\sum_{i=1}^{n} X_{i} - \mu\right)
      \cvL \mathcal{N}(0,\sigma^{2}).
  \]
\end{theorem}

\begin{remark}[The hypotheses are not interchangeable either]
  The loi forte asks for an espérance, that is $\expec(|X_{1}|) < +\infty$, and
  \emph{not} for a finite variance; the finite variance belongs to the weak law and
  to one proof of the loi forte, and importing it into the statement weakens the
  theorem for nothing. The théorème central limite does need a finite variance,
  since $\sigma^{2}$ appears in its conclusion. So the loi forte delivers the
  strongest mode under the weaker hypothesis, and the théorème central limite the
  weakest mode under the stronger one --- which is not a paradox: they say
  different things, one about where the mean goes and one about how large the
  error is on the way.
\end{remark}

\begin{conclusion}
Five symbols in these notes carry more than one meaning, because the course's
notation is kept as exams use it, collisions included, and each time the meaning is fixed by what
sits in the parentheses. $\mathcal{P}$ is the ensemble des parties
$\mathcal{P}(\Omega)$ and the loi de Poisson $\mathcal{P}(\lambda)$. $\mathcal{B}$
is the tribu de Borel $\mathcal{B}(\reals)$, the loi de Bernoulli
$\mathcal{B}(p)$, the loi binomiale $\mathcal{B}(n,p)$ and the loi Beta
$\mathcal{B}(a,b)$. $\lambda$ is the mesure de Lebesgue and the parameter of the
loi exponentielle and the loi de Poisson. $\Gamma$ is the loi Gamma
$\Gamma(a,\lambda)$ and the matrice de covariance of a vecteur gaussien,
$\mathcal{N}(\mu,\Gamma)$ --- and, in the densité of the loi Gamma, also the gamma
function. $\mathcal{E}$ is a generic tribu, as in $(E,\mathcal{E})$, and the loi
exponentielle $\mathcal{E}(\lambda)$. All five appear in the tables above in their
meaning as a loi, and two of them collide there on a single page: $\mathcal{B}$
carries the loi de Bernoulli, the loi binomiale and the loi Beta in the same two
tables, and the Gamma row prints the loi Gamma and the gamma function
side by side. That is why the notation column is the one to read first.
\end{conclusion}
